% Generated by `taskgen course examples` from arithmetic.dft_small (seed 3).
% Do not edit: change examples.yaml and regenerate.
\begin{example}[ДПФ над $\FF_p$ и теорема о свёртке]\label{ex:dft-f17}
В поле $\mathbb{F}_{17}$ дан элемент $\omega_0 = 4$. а)~Проверить, является ли $\omega_0$ главным корнем степени $n = 4$ из единицы; если нет, найти такой корень $\omega$ в виде $g^{(p-1)/n}$, где $g$ --- первообразный корень. б)~Вычислить ДПФ $X = \mathrm{DFT}_\omega(x)$, $Y = \mathrm{DFT}_\omega(y)$ для $x = (3, 3, 3, 0)$, $y = (7, 3, 0, 0)$, поточечное произведение $Z = XY$ и обратное ДПФ $z = \mathrm{DFT}^{-1}_\omega(Z)$. в)~Проверить ответ по теореме о свёртке, вычислив циклическую свёртку $x \circledast y$ непосредственно. Совпадает ли $z$ с линейной свёрткой $x * y$ над $\mathbb{Z}$?

\par\noindent\textit{Решение.} а) $p - 1 = 16$ делится на $n = 4$. Так как $n$ --- степень двойки и $n$ обратимо в $\mathbb{F}_{17}$, $\omega$ --- главный корень степени $n$ тогда и только тогда, когда $\omega^{n/2} = -1$ (в поле это равносильно тому, что порядок $\omega$ равен ровно $n$). Степени $\omega_0$: $\omega_0^{1} \equiv 4,\ \omega_0^{2} \equiv 16,\ \omega_0^{3} \equiv 13,\ \omega_0^{4} \equiv 1$. $\omega_0^{2} \equiv 16 = -1$, значит $\omega = \omega_0$ --- главный корень (порядок $4$).\par Нужные константы (по модулю $17$): $\omega^{-1} = \omega^{3} \equiv 13$, $n^{-1} \equiv 13$ ($4 \cdot 13 = 52 \equiv 1$). Степени $\omega^m$ (показатель $jk$ берётся по модулю $n$):\[\begin{array}{c|rrrr}m & 0 & 1 & 2 & 3 \\ \hline \omega^m & 1 & 4 & 16 & 13\end{array}\]б) $X_k = \sum_j x_j \omega^{jk}$, $Y_k = \sum_j y_j \omega^{jk}$ (нулевые слагаемые опущены):\[\begin{aligned}X_{0} &= 3 \cdot 1 + 3 \cdot 1 + 3 \cdot 1 = 9, \\ X_{1} &= 3 \cdot 1 + 3 \cdot 4 + 3 \cdot 16 = 63 \equiv 12, \\ X_{2} &= 3 \cdot 1 + 3 \cdot 16 + 3 \cdot 1 = 54 \equiv 3, \\ X_{3} &= 3 \cdot 1 + 3 \cdot 13 + 3 \cdot 16 = 90 \equiv 5, \\ Y_{0} &= 7 \cdot 1 + 3 \cdot 1 = 10, \\ Y_{1} &= 7 \cdot 1 + 3 \cdot 4 = 19 \equiv 2, \\ Y_{2} &= 7 \cdot 1 + 3 \cdot 16 = 55 \equiv 4, \\ Y_{3} &= 7 \cdot 1 + 3 \cdot 13 = 46 \equiv 12.\end{aligned}\]Поточечное произведение $Z_k = X_k Y_k \bmod 17$:\[\begin{array}{c|rrrr}k & 0 & 1 & 2 & 3 \\ \hline X_k & 9 & 12 & 3 & 5 \\ Y_k & 10 & 2 & 4 & 12 \\ \hline Z_k & 5 & 7 & 12 & 9\end{array}\]Обратное ДПФ: $u_j = \sum_k Z_k \omega^{-jk}$, $\omega^{-jk} = \omega^{(-jk) \bmod n}$, $z_j = n^{-1} u_j = 13 \cdot u_j$:\[\begin{aligned}u_{0} &= 5 \cdot 1 + 7 \cdot 1 + 12 \cdot 1 + 9 \cdot 1 = 33 \equiv 16, \\ u_{1} &= 5 \cdot 1 + 7 \cdot 13 + 12 \cdot 16 + 9 \cdot 4 = 324 \equiv 1, \\ u_{2} &= 5 \cdot 1 + 7 \cdot 16 + 12 \cdot 1 + 9 \cdot 16 = 273 \equiv 1, \\ u_{3} &= 5 \cdot 1 + 7 \cdot 4 + 12 \cdot 16 + 9 \cdot 13 = 342 \equiv 2;\end{aligned}\]\[\begin{array}{c|rrrr}j & 0 & 1 & 2 & 3 \\ \hline u_j & 16 & 1 & 1 & 2 \\ z_j = 13 u_j & 4 & 13 & 13 & 9\end{array}\]Итак, $z = (4, 13, 13, 9)$.\par в) Непосредственно: $(x \circledast y)_k = \sum_{i + j \equiv k\ (\mathrm{mod}\ 4)} x_i y_j$. Линейная свёртка $\ell = x * y = (21, 30, 30, 9)$:\[x \circledast y = (21, 30, 30, 9) \equiv (4, 13, 13, 9) \pmod{17},\]что совпадает с $z$ (теорема о свёртке). Однако $z$ \emph{не совпадает} с линейной свёрткой над $\mathbb{Z}$: коэффициент $30 \ge p$ известен лишь по модулю $p$. Чтобы получить линейную свёртку, нужны $n \ge 4$ и $p > \max_k (x * y)_k$.
\end{example}
